% Síntesis del banco aislado; resultados preservados y figuras retiradas del documento.
\Needspace{5\baselineskip}
\subsection{Validación funcional del autómata de seguridad en Stateflow}
\label{sec:validacion_seguridad_stateflow}
Se evaluó una copia del subsistema de seguridad mediante \textbf{25 escenarios de prueba seleccionados} para evaluar la respuesta del autómata de seguridad ante límites de posición, sobrevelocidad, parada de emergencia, pérdida de comunicación y solicitudes de reinicio. También se estudiaron fallas sucesivas, fallas simultáneas y condiciones peligrosas presentes desde el arranque.

Las pruebas se realizaron sobre el subsistema de seguridad aislado, mediante entradas programadas y con un período de ejecución de $20~\mathrm{ms}$. Se verificaron los permisos de movimiento y las autorizaciones de apertura de los frenos, sin incluir la respuesta mecánica de la planta. La \cref{tab:resumen_pruebas_seguridad} resume los principales resultados.

\begin{table}[!htb]
	\centering
	\small
	\renewcommand{\arraystretch}{1.2}
	\caption{Resumen de los resultados de las pruebas del autómata de seguridad.}
	\label{tab:resumen_pruebas_seguridad}
	\begin{tabular}{L{0.26\textwidth}L{0.66\textwidth}}
		\toprule
		\textbf{Aspecto evaluado} & \textbf{Resultado observado} \\
		\midrule
		
		Límites y sobrevelocidad &
		Los sensores se activaron al alcanzar el umbral, incluida la igualdad. \\
		
		Retención y reinicio &
		La emergencia permaneció enclavada después de retirar la causa. Un nuevo flanco de reinicio permitió recuperar los permisos cuando se cumplieron las condiciones correspondientes. \\
		
		Fallas sucesivas en ambos ejes &
		La segunda falla provocó el escalamiento a emergencia total. \\
		
		Fallas simultáneas en ambos ejes &
		Prevaleció la emergencia de izaje y el carro conservó su permiso. La prioridad de las transiciones y la detección por flanco impidieron reconocer posteriormente la falla horizontal. \\
		
		Pérdida del latido de comunicación &
		El watchdog detectó el vencimiento de los $5~\mathrm{s}$ sin cambios del latido. La emergencia total se activó $20~\mathrm{ms}$ después, debido al orden de ejecución de las regiones. \\
		
		Recuperación del watchdog &
		Fueron necesarios dos flancos de reinicio: uno para borrar la falla y otro para liberar la emergencia. Sin restablecer el latido, el fallo volvió a aparecer cinco segundos después. \\
		
		Fallas presentes al arrancar &
		Los límites y la sobrevelocidad inicialmente activos no provocaron intervención, porque la lógica esperaba un nuevo flanco. El pulsador, evaluado por nivel, sí produjo emergencia total. \\
		
		\bottomrule
	\end{tabular}
\end{table}

Los ensayos confirmaron la respuesta prevista ante las condiciones habituales y permitieron identificar limitaciones en la detección de fallas simultáneas, la protección durante el arranque y la recuperación del watchdog.

\par